\section{四维主角：$\mathcal{N}=4$ 超对称 Yang--Mills 与 S-对偶}
\label{sec:physics}

\begin{goals}
\begin{enumerate}
\item 记住 $\mathcal N=4$ SYM 的场内容与十六个超荷，知道它为什么是四维最大的无引力规范理论。
\item 会写复耦合 $\tau$，能一句话说清 S-对偶到底断言了什么。
\item 会读 Table 1，能解释 GNO 对偶群为什么正是 $\Ldual$，并分清这一讲里哪些是定理、哪些只是证据。
\end{enumerate}
\end{goals}

\subsection{从十维降下来：场内容与十六个超荷}

上一讲我们把四维量子场论当成一台黑箱用了：它在 $\Sigma\times C$ 上紧化，给出 target 为 $\MH$ 的二维 sigma model。这一讲把黑箱打开，顺便请出全篇第二个主角——S-对偶。

造 $\mathcal N=4$ SYM 最省事的办法是从十维降维。为什么偏偏是十维？Nahm 定理说，十维是超对称 Yang--Mills 能存在的最大维数；再往上，超荷作用会带进自旋 $3/2$ 的粒子，那已经是超引力 \src{B p.6}。十维的场只有两个：规范场 $A$（$G$-丛上的联络）与一个正手征旋量 $\lambda$（$S_+\otimes\mathrm{ad}(E)$ 的截面）\src{B p.7}。约定上 $A$ 与 $\lambda$ 取实李代数、在酉表示中反厄米，即 $A=-\ii A'$，与通常物理约定差一个 $\ii$，这样写是为了迁就数学文献 \src{B p.7 eq (2.2)}。$\Tr$ 是不变负定二次型，归一化到 $(1/8\pi^2)\int_{S^4}\Tr F\wedge F$ 取任意整数 \src{B p.8}。

十维的超对称由十六个守恒超荷 $Q_a$（$a=1,\dots,16$）实现，它们取值在 $S_-$ 里，满足原文 (2.9)：
\[
\{Q_a,Q_b\}=\sum_{I=1}^{10}\Gamma^I_{ab}\,P_I ,
\]
并且 $Q$ 与动量 $P_I$ 交换、按 $S_-$ 变换 \src{B p.9}。降维的手续是让所有场都不依赖 $x^4,\dots,x^9$。于是 $A$ 的四个分量给四维规范场 $A_\mu$，剩下六个分量给六个实标量 $\phi_i=A_{i+4}$（$i=0,\dots,5$）。十维曲率按指标落在哪个方向，自动分裂成 $F_{\mu\nu}$、$D_\mu\phi_i$、$[\phi_i,\phi_j]$ 三类，玻色部分就长成
\[
I_4=\frac{1}{e^2}\int d^4x\ \Tr\Bigl(\tfrac12 F_{\mu\nu}F^{\mu\nu}
+\sum_i D_\mu\phi_i D^\mu\phi_i+\tfrac12\sum_{i,j}[\phi_i,\phi_j]^2\Bigr)
\]
（原文 (2.10)，\src{B p.9}）。三块各有读法：规范场动能；六个标量的动能；以及标量之间的四次势——最后一项是正定平方和，这正是 $\mathcal N=4$ 有平坦势方向、能一直保超对称的原因。若再要求拉氏量在曲率上是二次的，那么给定 $G$ 之后它唯一到几个耦合常数，没有别的选择 \src{B p.6}。

费米子有几个？十维的 $\lambda$ 手征，降维后在 $Spin(1,3)\times Spin(6)$ 下分解。前一个因子是四维洛伦兹群，后一个因子满足 $Spin(6)\cong SU(4)$，这就是 \term{R-symmetry}{R-对称} 群 $SU(4)_R$。手征条件把十六维旋量表示按
\[
(2,1,4)\oplus(1,2,\bar 4)
\]
分解（原文 (2.14)，\src{B p.10}）。两块互为复共轭，加起来实维数仍是十六；四维超代数于是写成原文 (2.15)：
\[
\{Q^{AX},Q^Y_{\dot A}\}=\delta^X_Y\sum_{\mu=0}^{3}\Gamma^\mu_{A\dot A}P_\mu,
\qquad \{Q,Q\}=\{\bar Q,\bar Q\}=0,
\]
其中 $A,\dot A$ 是四维旋量指标，$X,Y=1,\dots,4$ 是 $SU(4)_R$ 指标，超荷总数 $4\times 4=16$ \src{B p.11}。同一个分解也作用在 $\lambda$ 上，所以费米子是四个 Weyl 费米子。场内容这一行就记全了：一个规范场 $A_\mu$、四个 Weyl 费米子、六个实标量。

顺便说一句：如果你觉得 1+4+6 这个组合眼熟，那不是巧合——它是十维 $10=1+4+6$ 的指标分裂，也是 $SU(4)_R$ 的 $\mathbf 4$ 加 $\bar{\mathbf 4}$ 拼出六个实标量的原因：$\mathbf 6=\mathbf 4\wedge\mathbf 4$。

\subsection{复耦合 $\tau$ 与 S-对偶}

$\mathcal N=4$ SYM 有两个经典耦合：规范耦合 $e$ 与 $\theta$ 角。$\theta$ 项是拓扑项，等于 $\theta$ 乘以第二陈数，也就是瞬子数 \src{B p.10}。量子力学只看作用量模 $2\pi$ 的整数倍，而 $\frac{1}{8\pi^2}\int\Tr F\wedge F$ 是整数，所以 $\theta$ 是周期变量 $\theta\sim\theta+2\pi$。把两者合成一个复数
\[
\tau=\frac{\theta}{2\pi}+\frac{4\pi\ii}{e^2},
\]
则因 $\Im\tau>0$，$\tau$ 落在上半平面 $\HH$ 里（原文 (2.12)，\src{B p.10}）。$SL(2,\RR)$ 按 $\tau\mapsto(a\tau+b)/(c\tau+d)$、$ad-bc=1$ 作用。其中 $T:\tau\mapsto\tau+1$ 只是把 $\theta$ 平移 $2\pi$，在任何 $G$ 下都是经典对称 \src{B p.11}。

真正的重头是 $S$。原文的一句话是：\ldots an additional quantum symmetry that inverts $\tau$, exchanges $G$ with $\Ldual$, and exchanges electric and magnetic charges \src{B p.11}。请盯住三个词：\emph{additional}（它不是已有的经典对称）、\emph{quantum}（只在量子理论里成立）、\emph{inverts}（把耦合倒过来）。$S$ 与 $T$ 一起生成 $SL(2,\RR)$ 的一个离散子群 $\Gamma$。$G$ 属 simply-laced 型时，$S$ 就是原文 (2.16)：
\[
S=\begin{pmatrix}0&1\\-1&0\end{pmatrix},\qquad
\tau\mapsto-\frac{1}{\tau},
\]
与 $T$ 生成 $SL(2,\ZZ)$ \src{B p.11}。$G$ 非 simply-laced 时要稍作调整：
\[
\tau\mapsto-\frac{1}{n_g\tau},\qquad
n_g=(|\alpha_{\mathrm{long}}|/|\alpha_{\mathrm{short}}|)^2 ,
\]
即 $n_g=2$（$F_4$）、$n_g=3$（$G_2$），由原文 (2.17) 的矩阵实现 \src{B p.11}。$\tau$ 写成 $4\pi\ii/e^2$ 而不是 $\ii/e^2$ 的归一化，正好让 simply-laced 情形下 $S$ 是干净的对合。

S 到底交换了什么？图~\ref{fig:sdual} 把这句话画成一张图：$\tau$ 与 $G$ 一起被换掉，而电与磁的角色互换。图上写了 $\tau'=-1/(n_g\tau)$，意思是右边的耦合是左边耦合的变换值；右边那个 $\Ldual$ 并不是同一个群，它是下一节的主角。

\begin{figure}[htbp]
\centering
\begin{tikzpicture}[>=Stealth, font=\small,
  side/.style={draw=kwblue,thick,rounded corners=3pt,fill=kwblue!5,
    align=center,inner sep=7pt,text width=4.1cm}]
  \node[side] (g) {$G$ 侧\\[3pt]
    $\tau=\theta/2\pi+4\pi\ii/e^{2}$\\[3pt]
    电激发：$G$ 的权格};
  \node[side,right=2.2cm of g] (dg) {$\Ldual$ 侧\\[3pt]
    $\tau'=-1/(n_g\tau)$\\[3pt]
    磁激发：$\Ldual$ 的权格};
  \draw[<->,very thick,kwred] (g.east) -- node[above=2pt,font=\small]{$S$}
    node[below=2pt,font=\scriptsize,align=center]{电 $\leftrightarrow$ 磁\\$G\leftrightarrow\Ldual$} (dg.west);
\end{tikzpicture}
\caption{S-对偶交换的三件事：把 $\tau$ 换成 $-1/(n_g\tau)$、把 $G$ 换成 $\Ldual$、把电荷与磁荷互换。它是量子对称，不在经典拉氏量里。}
\label{fig:sdual}
\end{figure}

还有一个细节第 3 讲会直接用到：$S$ 在超荷上不是平凡地作用。原文特别说明，这件事不能从经典拉氏量读出来，因为除了 $T$ 生成的子群以外，$\Gamma$ 根本不含经典对称 \src{B p.13}。量子推理给出的答案是：$S$ 在 $\mathbf 4$（即 $Q^{AX}$）上乘 $\ee^{\ii\hat\phi}$、在 $\bar{\mathbf 4}$ 上乘 $\ee^{-\ii\hat\phi}$，而
\[
\ee^{-\ii\hat\phi(\gamma)}=\Bigl(\frac{|c\tau+d|}{c\tau+d}\Bigr)^{1/2}
\]
（原文 (2.25)，\src{B p.15}）；平方根意味着作用在超荷上的是 $\Gamma$ 的双覆盖，要由 $(-1)^F$ 扩一下才对得上。为什么这件事要紧？因为第 3 讲扭化所选的，正是“哪些超荷在扭化后还活着”，所以 $S$ 怎样作用在超荷上，直接决定 S-对偶在扭化理论里变成什么。这一句可以先记下，不必现在算。

\subsection{GNO、Montonen--Olive，以及为什么是 $\Ldual$}

S-对偶不是凭空冒出来的，它有两个来源，而这两个来源正好构成本讲的骨架。

第一个是 GNO（Goddard--Nuyts--Olive，1976）。他们证明：在规范理论里，电荷取规范群的权格，而磁荷取“对偶群”的权格 \src{B p.2}；一般紧李群的磁荷分析更早由 Englert--Windey 给出 \src{B p.2}。为什么磁荷会落到别处？因为 't~Hooft--Polyakov 单极子不是微扰态：它的磁荷由解在无穷远处的绕数决定，而绕数的取值取决于 $G$ 里那个 $U(1)$ 子群怎么嵌入，也就是余权格。于是“能承载这种磁荷”的群，自然是 $G$ 的对偶群。原文在 §2 末的 Coulomb 分支讨论里把这件事写得更细：电、磁荷分别取权格与余权格，中心荷依赖 $\tau$ \src{B pp.14--15}。

第二个是 Montonen--Olive（1977）。他们比 GNO 更进一步，提出一个猜想：以 $G$ 为基础的那个规范理论，与把 $e$ 换成 $1/e$、把群换成对偶群的那个理论，是同一个理论 \src{B p.2}。GNO 说的是“磁荷住在哪个格”，MO 说的是“整个理论等价”；后者强得多，因为理论等价不仅要求谱一致，还要求所有关联函数一致。

两条线在 1977 年底合流。原文写：GNO 对偶群 actually the same as the Langlands dual group $\Ldual$，这是 Atiyah 观察到的，他当时建议第二作者注意 Langlands 纲领与量子场论的联系，并推荐了 GNO 与 MO 两篇文章 \src{B p.3}。之后是三步：超对称规范理论里 MO 对偶更自然 \src{B p.3}；$\mathcal N=4$ 是它最合适的场所；$\theta$ 角把原来的 $\ZZ_2$ 对偶扩成 $SL(2,\ZZ)$。

为什么 GNO 对偶群正是 $\Ldual$？本讲只给格层面的解释，因为原文本身也只给了这句话。Langlands 对偶在群论上就是把根与余根互换，于是权格与余权格互换；而磁荷量子化恰恰要求磁荷落在 $G$ 的余权格上，而“$G$ 的余权格”与“$\Ldual$ 的权格”是同一组格。就这一步，但值得写下来。单连通性的差别则表现为中心商：$SU(N)$ 与 $PSU(N)=SU(N)/\ZZ_N$ 的中心与基本群互换，这正是表~\ref{tab:dual} 第二行。再补一句诚实提醒：$\pi_1(G)\neq0$ 时还要看特征类 $\xi\in H^2(M,\pi_1(G))$，不同 $\xi$ 的配分函数按对偶群的酉表示变换 \src{B p.12}，这里不详谈。

\subsection{Table 1 怎么读，以及为什么这只是“证据”}

Table 1 只有两列：左边 $G$，右边 $\Ldual$。原文的七行照抄如下（只去掉排版）\src{B p.3}：

\begin{table}[htbp]
\centering
\small
\begin{tabular}{@{}ll@{}}
\toprule
$G$ & $\Ldual$ \\
\midrule
$U(N)$ & $U(N)$ \\
$SU(N)$ & $PSU(N)=SU(N)/\ZZ_N$ \\
$Spin(2n)$ & $SO(2n)/\ZZ_2$ \\
$Sp(n)$ & $SO(2n+1)$ \\
$Spin(2n+1)$ & $Sp(n)/\ZZ_2$ \\
$G_2$ & $G_2$ \\
$E_8$ & $E_8$ \\
\bottomrule
\end{tabular}
\caption{Table 1 的简化复制：$G$ 与其 Langlands（GNO）对偶群 $\Ldual$，见 \src{B p.3}。}
\label{tab:dual}
\end{table}

怎么读？三条读法。

第一，右列不是“随便另一个群”，而是由 $G$ 的根数据对偶地定出来的。$G$ 单连通时 $\Ldual$ 是伴随型，所以右列常出现商：$SU(N)$ 对偶到 $PSU(N)$，$Sp(n)$ 对偶到 $SO(2n+1)$，$Spin(2n)$ 对偶到 $SO(2n)$ 再除中心。对偶是对合的：把表倒过来读，$PSU(N)$ 的对偶正是 $SU(N)$。

第二，中心与基本群互换。$SU(N)$ 的中心是 $\ZZ_N$，$PSU(N)$ 的基本群是 $\ZZ_N$；这一行是“磁荷取余权格”最干净的体现。同理，$Sp(n)$ 那一行的奥妙在于：$Sp(n)$ 与 $Spin(2n+1)$ 是仅有的“李代数不同构”的一对对偶群，其余各行两边都是同一个李代数；正因为李代数不同构，$\tau$ 在两边的归一化要重新对齐，$n_g=2$ 就是这么来的。

第三，自对偶行要小心。$G_2$ 与 $E_8$ 的根与余根同构，所以自对偶；$U(N)$ 行也自对偶，但严格说 $U(N)$ 不是半单群，它的 $U(1)$ 因子要单独处理——§1 写这张表时还没有 §2 的“$G$ 紧”假设，读的时候要知道这个差别。还有一处缺席值得注意：表里没有 $F_4$。它作为群是自对偶的，但它非 simply-laced，$S$ 作用在 $\tau$ 上是 $\tau\mapsto-1/(2\tau)$。这个小小的缺席恰好说明：群的对偶与耦合的对偶是两个层次，前者是群论，后者还要看归一化。

最后把这一讲的地位说清楚：以上全是期望与证据，不是定理。原文自己用的词就是 conjecture 与 asserts——$S$ 是被“断言”为一个量子对称的 \src{B p.11}。四十多年来它有一大批互相独立的证据：BPS 谱的匹配、瞬子配分函数的计算、弦论里的实现，以及整本 \cite{kw06} 的前提——如果 S-对偶不成立，几何 Langlands 的物理推导根本起不来。但“证据极多”不等于“已经证明”。到 2026 年为止，四维 $\mathcal N=4$ SYM 的 S-对偶仍然没有数学证明。这也解释了 K--W 的写法：它把 S-对偶当作公理式的输入，推出一整套几何 Langlands 的陈述；而真正给出独立证明的，是 2024 年前后 Gaitsgory 等人的纯代数几何工作，那条路线完全不经过四维场论。物理路线与数学证明是两条独立的路，最后指向同一个答案——这是全篇最让人服气的地方，也是最需要你自己判断的地方。

\begin{keybox}
$\mathcal N=4$ SYM 是 S-对偶最干净的场所：$S$ 把 $\tau$ 换成 $-1/(n_g\tau)$、把 $G$ 换成 $\Ldual$、把电与磁互换。$\Ldual$ 正是 GNO 1976 从磁荷量子化读出的对偶群，也是 Langlands 对偶群；但这一整套是猜想与证据，不是定理。
\end{keybox}

\begin{follow}
\begin{itemize}
\item \textbf{B §1，pp.2--5}：GNO、Montonen--Olive、Atiyah 的观察、Table 1 与后续发展。全篇所有“为什么”的来源，先读它再读技术细节。
\item \textbf{B §2.1，pp.6--11}：场内容、十六超荷、降维、$\tau$ 与 $SU(4)_R$。沿 (2.2)--(2.12) 读，盯住反厄米约定与 $\Tr$ 的归一化。
\item \textbf{B §2.2，pp.11--16}：S-对偶的断言、(2.16)/(2.17)，以及 pp.13--15 超荷上的相位 $\hat\phi$。这一段是第 3 讲扭化的直接前置。
\end{itemize}
\end{follow}

\begin{checkbox}
\begin{itemize}
\item 写出 $\mathcal N=4$ SYM 的场内容：规范场、实标量、Weyl 费米子各几个？它们分别从十维的哪个场来？
\item $\theta$ 为什么是周期变量？$\tau$ 为什么落在上半平面？为什么 $T:\tau\mapsto\tau+1$ 是经典对称而 $S$ 不是？
\item S-对偶断言的三件事是什么？非 simply-laced 时 $\tau\mapsto-1/(n_g\tau)$，$n_g$ 是什么，为什么 $F_4$ 对应 $n_g=2$？
\item GNO 1976 说了什么？为什么单极子的磁荷落在余权格而不是权格？
\item 读 Table 1：$SU(N)$ 一行的中心与基本群怎样互换？为什么表里没有 $F_4$？
\item 这一讲里哪些是定理、哪些是猜想？S-对偶对 $\mathcal N=4$ 被证明了吗？若没有，K--W 凭什么用它？
\end{itemize}
\end{checkbox}
